% ============================================================
% appendix.tex —— 附录 A–D
% ============================================================
\appendix

\chapter{土名对照表：你发明的，大学的叫法}
\label{app:tuming}

\begin{center}
\begin{longtable}{p{0.31\textwidth} p{0.31\textwidth} c}
\toprule
本书土名 & 正式术语 & 首次出场 \\
\midrule
\endhead
余数王国、钟表世界、压扁 & 模 $n$ 算术、同余 & 第 1 章 \\
裁边角 & 欧几里得算法（辗转相除） & 第 2 章 \\
原子的配方 & 素因数分解（唯一分解定理） & 第 3 章 \\
假设敌人存在 & 反证法 & 第 4 章 \\
各管一摊 & 中国剩余定理 & 第 5 章 \\
$1$ 之外全装死 & 互素、裴蜀等式 & 第 2、5 章 \\
体检四项 & 群的公理 & 第 7 章 \\
什么都不做的数 & 单位元 & 第 7 章 \\
搭档 & 逆元 & 第 7 章 \\
复原动作全家福 & 对称群（二面体群） & 第 8 章 \\
失散的双胞胎 & 同构 & 第 8 章 \\
群中之群 & 子群 & 第 9 章 \\
归零周期 & 元素的阶 & 第 9 章 \\
模具扣出的盘子 & 陪集 & 第 9 章 \\
并肩干活的世界 & 环 & 第 10 章 \\
除法全面通行 & 域 & 第 10 章 \\
假零制造者 & 零因子 & 第 10 章 \\
分数流水线 & 有理数域的构造（分式域） & 第 10 章 \\
另一种数 & 多项式环 & 第 11 章 \\
塔内居民 & 可作图数（域的二次扩张塔） & 第 12 章 \\
落脚点 & 极限 & 第 14 章 \\
要多近有多近的比赛 & $\varepsilon$--$N$ 定义 & 第 14 章 \\
无限连加 & 无穷级数 & 第 13 章 \\
瞬间变化率 & 导数 & 第 15 章 \\
求导机器 & 微分（求导） & 第 15 章 \\
穿盒术 & 链式法则 & 第 16 章 \\
堆积、切条 & 定积分、黎曼和 & 第 17 章 \\
反导数、后门 & 原函数、牛顿--莱布尼茨公式 & 第 18 章 \\
指数的反问题 & 自然对数 $\ln$ & 第 19 章 \\
素数人口 & 素数计数函数 $\pi(x)$ & 第 20 章 \\
越晚越贴 & 渐近等价 & 第 20 章 \\
分块放大术 & 比较判别法（前身） & 第 21 章 \\
齿轮、翻译通道 & 欧拉乘积、几何级数 & 第 22 章 \\
溺水点 & 零点 & 第 24 章 \\
藏宝图 & $\zeta$ 的解析延拓与黎曼猜想 & 第 24 章 \\
弯度、转弯强度 & 曲率 $\kappa$ & 第 25 章 \\
贴身圆 & 密切圆 & 第 25 章 \\
绷紧的线 & 测地线 & 第 26 章 \\
世界的弯曲度 & 高斯曲率 $K$ & 第 27 章 \\
圆圈缩水/虚胖 & 正/负曲率的圆周实验 & 第 27 章 \\
画家平面 & 射影平面 & 第 30、31 章 \\
交点乘法定律 & 贝祖定理 & 第 30 章 \\
发身份证 & 齐次坐标 & 第 31 章 \\
循环点 & 圆的公共无穷远点 $[1:\pm\mathrm{i}:0]$ & 第 31 章 \\
第三点传送门、弦切游戏 & 椭圆曲线的群律 & 第 32 章 \\
\bottomrule
\end{longtable}
\end{center}

\chapter{符号表}

\begin{center}
\begin{longtable}{lp{0.72\textwidth}}
\toprule
符号 & 含义（首次出场） \\
\midrule
\endhead
$a \equiv b \pmod n$ & $a$ 与 $b$ 除以 $n$ 余数相同（第 1 章） \\
$\gcd(a,b)$ & $a,b$ 的最大公因数（第 2 章） \\
$\heartsuit,\ \clubsuit$ & 本书自造运算（序章） \\
$n!$ & 阶乘 $1\times2\times\cdots\times n$（第 4 章） \\
$\lim$ & 落脚点（极限，第 14 章） \\
$f'(x)$ / $\dfrac{ds}{dt}$ & 导数（第 15 章） \\
$\displaystyle\int_a^b f\,dt$ & 从 $a$ 到 $b$ 的定积分（第 17 章） \\
$\mathrm{e},\ \ln$ & 自然常数与自然对数（第 19 章） \\
$\pi(x)$ & $x$ 以内素数个数（第 20 章） \\
$\displaystyle\sum,\ \prod$ & 连加、连乘（第 13、22 章） \\
$\zeta(s)$ & $\displaystyle\sum_{n\ge1} n^{-s}$（第 24 章） \\
$\kappa,\ K$ & 曲率、高斯曲率（第 25、27 章） \\
$[X:Y:Z]$ & 齐次坐标（第 31 章） \\
$O$ & 椭圆曲线群的无穷远点（第 32 章） \\
$\star\ /\ \bigstar$ & 习题难度标记 \\
\bottomrule
\end{longtable}
\end{center}

\chapter{初中知识使用地图}

\begin{center}
\begin{longtable}{p{0.34\textwidth} p{0.55\textwidth}}
\toprule
初中知识点 & 被回炉重造的位置 \\
\midrule
\endhead
除法竖式与余数 & 第 1 章（余数王国）、第 13 章（循环小数）、第 17 章（切条求和） \\
公因数与公倍数 & 第 2、3 章（铺地砖、原子配方）、第 10 章（除法通行证） \\
乘法九九表 & 第 8 章（群表）、第 10 章（逆元表） \\
乘法公式 $(a\pm b)^2$ & 第 15 章（发明导数的核心机器） \\
一元二次方程与判别式 & 第 16 章（顶点公式）、第 29--30 章（切点两票） \\
因式分解 & 第 11 章（因式定理）、第 30 章（三次方程三根） \\
平面直角坐标系 & 第 29 章（方程动物园）、第 31 章（齐次坐标） \\
一次与二次函数图像 & 第 15 章（切线）、第 29 章（交点 $=$ 联立） \\
三角形内角和、圆的性质 & 第 26 章（球面 $270^\circ$）、第 31 章（圆的共同秘密） \\
平均速度 $\times$ 时间 & 第 15 章（瞬时速度）、第 17 章（里程 $=$ 速度堆积） \\
统计初步（占比、普查） & 第 20 章（素数人口普查）、第 23 章（误差分析） \\
正负数运算 & 第 1 章（余数房间）、第 6 章（$31 \equiv -2$ 加密解密） \\
\bottomrule
\end{longtable}
\end{center}

\chapter{下一站书单：当你准备好了}

\begin{description}[leftmargin=2em]
  \item[通识再进一步] R. 柯朗、H. 罗宾《什么是数学》——本书的"成年版"，六篇全部有对应章节；J. Wilkes《烧掉数学书》——本书的精神源头，英文原版 \textit{Burn Math Class}。
  \item[分析（篇三的完全体）] T. 陶哲轩《陶哲轩实分析》——从自然数公理出发严格重建极限；华罗庚《高等数学引论》——中国人的微积分开篇。
  \item[代数（篇二的完全体）] M. Artin《代数》——群环域与伽罗瓦理论的第一站；丘维声《抽象代数基础》——中文首选。
  \item[数论（篇一的完全体）] G. H. Hardy、E. M. Wright《数论导引》——经典中的经典；潘承洞、潘承彪《初等数论》——中文双璧。
  \item[解析数论（篇四）] T. Apostol《解析数论导引》——第 23、24 章地毯下的全部机器；J. Derbyshire《素数之恋》——黎曼猜想的最佳观光导游。
  \item[微分几何（篇五）] M. P. do Carmo《曲线与曲面的微分几何》——标准教材；陈省身、陈维桓《微分几何讲义》——大师的中文讲义。
  \item[代数几何与椭圆曲线（篇六）] J. Silverman、J. Tate《Rational Points on Elliptic Curves》——第 32 章之后的下一章；M. Reid《Undergraduate Algebraic Geometry》——最轻的入门。
\end{description}

\vspace{2em}
\noindent 祝下一程顺利。记住本书第一页的许可证：\textbf{它永久有效}。
